%% kapa: exam=2 round=21 year=2010 period=1 subject=appraisal_practice subject_ko=감정평가실무 seq=4 number=4 points=15
%% source: Q-Net 감정평가사 기출문제 (cst00309) artlSeq=5024183 file=extracted/1008/5024183/2010년 제21회 감정평가사 2차시험 기출문제/1교시.hwp
%% conversion: kordoc 4.24.51 -> tools/gen2.py (auto); review=proofread-sonnet; flags=none
\begin{longquestion}{【문제 4】}
이대한씨는 甲구 乙동에 소재한 구형 단독주택을 구입하여 철거한 후, 다가구주택을 신축하여 임대하고자 한다. 다음의 자료를 활용하여 신축예정 다가구주택의 수익가치를 산정하고, 이 사업으로 인한 요구수익률을 충족시킬 수 있는 구형 단독주택의 최대 매수(지불)가능가격을 산정하시오. (15점) \par
<자료 1> 신축예정 다가구주택의 개요 \par
1) 총 18개의 단위호로 구성된 철근콘크리트조의 원룸형으로 기존주택 철거비를 포함하여 신축공사비는 총 550,000,000원이 소요될 예정임. \par
2) 신축 건물의 예상 임대료는 각 단위호당 임대보증금 10,000,000원에 월세 800,000원을 받을 수 있을 것으로 추정되며, 월세 징구에 따른 시간가치는 고려하지 아니함. \par
3) 임대보증금의 운용이율은 1년만기 정기예금금리 수준인 연 4\%가 적정한 것으로 조사되며, 예상되는 공실손실상당액 및 대손액 등은 가능총수익(PGI)의 8\%, 관리 등에 소요되는 제반 경비는 유효총수익(EGI)의 20\%로 예상됨. \par
<자료 2> 인근 유사 다가구주택의 거래사례 및 임대자료 (사례 수 : 3개) \par
1) 사례1 : 2009년 1월 - 매매금액 2,000,000,000원 \par
연간 순수익(NOI) 220,000,000원 \par
2) 사례2 : 2009년 6월 - 매매금액 2,000,000,000원 \par
연간 순수익(NOI) 180,000,000원 \par
3) 사례3 : 2010년 8월 - 매매금액 1,800,000,000원 \par
연간 순수익(NOI) 180,000,000원 \par
<자료 3> 기타 참고사항 \par
1) 이대한씨는 사업자본의 50\%를 차입(타인자본)할 예정임. 차입자금의 연간 금리수준은 10\%이며, 자기자본에 대한 요구수익률은 위험을 고려하여 차입자금 금리수준의 2배를 기대하고 있음. 한편, 본 사업에 적용할 종합적인 요구수익률은 물리적 투자결합법을 활용하여 산정할 예정임. \par
2) 수익가치 산정에 사용할 환원이율은 시장추출법에 의하여 산정하되, 수집된 자료 중 가장 최근의 자료에 50\%의 가중치를 두고, 나머지 자료는 동일한 비중으로 취급할 예정이고, 환원방법은 직접법에 의함. \par
\end{longquestion}
